\chapter{Gamma-Glutamyl Transferase (GGT)}

\section*{What Is This Test?}
Gamma-glutamyl transferase (GGT) is a microsomal enzyme found on the surface of cells with high secretory or absorptive activity, primarily in the liver (canalicular membrane of hepatocytes and bile duct epithelium), kidneys (proximal tubule), pancreas, and intestine. GGT catalyzes the transfer of gamma-glutamyl groups from peptides (such as glutathione) to amino acids, playing a key role in glutathione metabolism and the transport of amino acids across cell membranes. In the liver, GGT is located on the canalicular membrane of hepatocytes and in the epithelium of bile ductules. Serum GGT is elevated in any condition causing cholestasis (impaired bile flow) because bile acids solubilize GGT from the canalicular membrane. GGT is also induced by drugs and alcohol that stimulate hepatic microsomal enzymes (cytochrome P450 system). This dual sensitivity---to cholestasis and to microsomal enzyme induction---makes GGT both extremely sensitive (elevated in many hepatobiliary conditions) and relatively non-specific. However, GGT's greatest clinical utility is as a confirmatory marker that elevated ALP is of hepatobiliary origin rather than bone: GGT is not present in bone, so a normal GGT in the setting of elevated ALP indicates a bone (or other non-hepatic) source. GGT is the most sensitive marker for detecting alcohol use and for monitoring abstinence---it rises with chronic alcohol consumption and falls with abstinence (half-life approximately 7--10 days) \cite{whitfield2001gamma}.

\section*{Alternative Names}
Gamma-GT; Gamma-Glutamyl Transpeptidase; GGTP; Serum GGT

\section*{Principle}
GGT is measured by a kinetic spectrophotometric method using gamma-glutamyl-3-carboxy-4-nitroanilide as the substrate. In the presence of glycylglycine (acceptor), GGT transfers the gamma-glutamyl group to glycylglycine, releasing 5-amino-2-nitrobenzoate, which absorbs light at 410 nm: L-gamma-glutamyl-3-carboxy-4-nitroanilide + glycylglycine $\rightarrow$ 5-amino-2-nitrobenzoate + gamma-glutamyl-glycylglycine. The rate of increase in absorbance at 410 nm is directly proportional to GGT activity. The IFCC standardized method was published in 2002 \cite{ifcc2002gamma}. The test is performed at 37 degrees C on automated chemistry analyzers. Results are reported in IU/L. This test is NOT performed by hematology analyzers \cite{tietz2023textbook}.

\section*{History}
GGT was first described by Hanes and colleagues in 1952 as a kidney enzyme. Its clinical utility in liver disease was established by Szczeklik and Or{\l}owski in the 1960s. The association of GGT with alcohol consumption was recognized in the 1970s, leading to its widespread use as a marker for alcohol abuse. Rosalki and Rau described the kinetic method for GGT measurement in 1972 \cite{rosalki1972serum}.

\section*{Why This Test Is Ordered}
\begin{itemize}
  \item Confirmation of hepatobiliary origin of elevated alkaline phosphatase: GGT or 5'-nucleotidase determines whether elevated ALP is from liver/bile ducts (both elevated) or bone (GGT normal)
  \item Screening for alcohol use disorder and monitoring of alcohol abstinence: GGT is the most sensitive blood marker for chronic alcohol consumption, elevated in 70--80\% of heavy drinkers \cite{conigrave2003traditional}
  \item Evaluation of cholestasis and biliary obstruction alongside ALP
  \item Detection of drug-induced hepatic microsomal enzyme induction: anticonvulsants (phenytoin, phenobarbital, carbamazepine), rifampin, warfarin, oral contraceptives
  \item Monitoring of alcohol abstinence in patients in treatment programs: GGT falls with a half-life of approximately 7--10 days after alcohol cessation
  \item Screening for bile duct injury/disease in conditions like primary biliary cholangitis, primary sclerosing cholangitis
  \item Evaluation of suspected metastatic liver disease: GGT is often the earliest and most markedly elevated liver enzyme in hepatic metastases
\end{itemize}

\section*{Primary vs Secondary Test}
GGT is a secondary, adjunctive test. It is not part of the routine BMP or CMP. Its primary use is to determine whether a patient with elevated ALP has hepatobiliary disease (elevated GGT) or bone disease (normal GGT). It is also used in alcohol screening and monitoring.

\section*{How This Test Is Performed}
A healthcare professional draws blood from a vein into a serum separator tube (gold-top) or lithium heparin tube (green-top). Standard venipuncture technique is used. The sample is centrifuged and GGT activity is measured on an automated chemistry analyzer. Results are available within 1--2 hours. Hemolyzed samples are acceptable for GGT measurement (unlike AST and LDH, which are markedly affected by hemolysis).

\section*{What Preparation Is Needed}
Fasting for 8 hours is recommended. Alcohol consumption within 24--48 hours before the test can transiently elevate GGT. The following medications induce GGT: phenytoin, phenobarbital, carbamazepine, valproic acid, rifampin, isoniazid, warfarin, oral contraceptives, and barbiturates. These should be noted because GGT elevation may represent hepatic enzyme induction rather than liver injury \cite{wallach2022interpretation}. Conversely, clofibrate and certain other fibrates can lower GGT. GGT may be modestly elevated in obesity and metabolic syndrome (likely from associated NAFLD). Smoking may slightly elevate GGT.

\section*{Contraindications and Precautions}
\begin{itemize}
  \item No absolute contraindications. Important considerations:
  \item (1) GGT is very sensitive but not specific. Common causes of GGT elevation that are not clinically significant include: recent alcohol consumption, many commonly prescribed medications (enzyme induction), obesity, and even moderate caffeine intake.
  \item (2) GGT may be more sensitive than ALP for detecting hepatic metastases, but it lacks specificity and should not be used as a screening test in the absence of other abnormal liver tests.
  \item (3) In acute viral hepatitis, GGT rises and falls alongside aminotransferases but is less useful than ALT for monitoring.
  \item (4) GGT is not useful for detecting cirrhosis or assessing synthetic function---patients with advanced cirrhosis may have normal GGT.
  \item (5) GGT may be elevated in chronic kidney disease (mechanism unclear, possibly from enzyme induction or decreased clearance).
  \item (6) GGT is normal in bone disease; a normal GGT confirms that elevated ALP is from bone, intestinal, or placental source.
  \item (7) A markedly elevated GGT with disproportionately low ALP suggests drug- or alcohol-induced enzyme induction without significant cholestasis.
  \item (8) GGT rises in pregnancy (especially third trimester) and in neonates.
  \item (9) Extremely high GGT (>10$\times$ ULN) with ALP >3$\times$ ULN strongly suggests biliary obstruction or diffuse hepatic infiltration.
\end{itemize}
\section*{Risks}
Minimal risk from venipuncture. Mild pain or bruising resolves in 1--2 days.

\section*{What Happens After the Test}
Results available within 1--2 hours. Interpretation depends on concurrent ALP and liver tests \cite{kwo2017acg}: (1) Elevated ALP + elevated GGT: hepatobiliary disease. Determine if cholestatic (ALP > ALT elevation), infiltrative, or drug-induced. (2) Elevated ALP + normal GGT: bone disease (Paget, fractures, bone metastases, vitamin D deficiency, growing child), intestinal ALP rise after fatty meal, or pregnancy. (3) Isolated GGT elevation (normal ALP, normal ALT/AST): hepatic microsomal enzyme induction from alcohol, anticonvulsants, or other drugs; NAFLD; or early cholestasis. Assess alcohol intake and medication list. (4) GGT >10$\times$ ULN with ALP >3$\times$ ULN: biliary obstruction (choledocholithiasis, malignancy, PSC, PBC). (5) Monitoring alcohol abstinence: GGT should fall by approximately 50\% every 7--10 days after alcohol cessation; normalization takes 2--6 weeks.

\section*{Understanding Results}

\subsection*{Below Normal}
\begin{itemize}
  \item Low GGT is not clinically significant and requires no evaluation.
  \item Normal GGT in a patient with elevated ALP strongly suggests bone, intestinal, or placental origin.
  \item Some individuals have constitutively low GGT (genetic variation). \cite{schiele1977gamma}
  \item Clofibrate and certain other drugs may lower GGT.
  \item Advanced cirrhosis may have normal GGT due to loss of hepatic enzyme-producing cells.
\end{itemize}

\subsection*{Above Normal}
\begin{itemize}
  \item \textbf{Hepatobiliary disease (with elevated ALP):} (a) Biliary obstruction: choledocholithiasis, pancreatic head mass, cholangiocarcinoma, PSC, PBC, benign strictures---GGT rises earlier and persists longer than ALP in biliary obstruction. (b) Infiltrative liver disease: hepatic metastases (GGT is often the earliest and most sensitive marker, elevated before ALP in some cases), lymphoma, sarcoidosis, granulomatous diseases. (c) Drug-induced cholestasis: estrogens, anabolic steroids, phenothiazines, erythromycin, chlorpromazine, amoxicillin-clavulanate, cyclosporine. (d) Cholecystitis (acute or chronic). (e) Viral hepatitis (GGT is elevated but less useful than ALT for monitoring).
  \item \textbf{Hepatic microsomal enzyme induction (with normal ALP):} (a) Alcohol use: GGT rises after 2--4 weeks of heavy alcohol consumption (>60 g/day) and is elevated in 70--80\% of heavy drinkers. GGT may be the only abnormal liver test in early alcoholic liver disease. (b) Medications: anticonvulsants (phenytoin > phenobarbital > carbamazepine > valproic acid), rifampin, isoniazid, warfarin, oral contraceptives, barbiturates, glutethimide. (c) Obesity and metabolic syndrome (likely associated NAFLD). (d) Diabetes mellitus.
  \item \textbf{Pancreatic disease:} Acute and chronic pancreatitis; pancreatic carcinoma. GGT is elevated in approximately 70--90\% of patients with pancreatic malignancy.
  \item \textbf{Others:} Chronic kidney disease, acute myocardial infarction (transient elevation in the days after MI), hyperthyroidism, COPD, and critical illness.
\end{itemize}

\section*{Normal Results}
\begin{itemize}
  \item \textbf{Adult males:} 8--61 IU/L (varies by laboratory; typical: 0--50 IU/L)
  \item \textbf{Adult females:} 5--36 IU/L (lower than males due to sex hormone differences)
  \item \textbf{Children:} 5--32 IU/L
  \item \textbf{Neonates:} Up to 5$\times$ adult ULN (high at birth, declines rapidly)
  \item \textbf{Pregnancy (third trimester):} Mildly elevated (1.5--2$\times$ ULN)
  \item \textbf{After alcohol abstinence:} Falls by ~50\% per 7--10 days
  \item \textbf{GGT half-life:} 7--10 days (up to 28 days in some studies)
\end{itemize}

\section*{Follow-up Tests}
\begin{itemize}
  \item \textbf{ALP, ALT, AST, bilirubin (total and direct), albumin, INR:} Complete liver panel to determine pattern and severity.
  \item \textbf{5'-Nucleotidase:} Alternative confirmatory test for hepatobiliary origin of ALP. Like GGT, it is not present in bone.
  \item \textbf{Abdominal ultrasound:} First-line imaging for biliary obstruction, gallstones, liver metastases.
  \item \textbf{MRCP or ERCP:} If biliary obstruction suspected.
  \item \textbf{Carbohydrate-deficient transferrin (CDT) and ethyl glucuronide (EtG):} More specific markers of chronic alcohol use than GGT. CDT has higher specificity.
  \item \textbf{Anti-mitochondrial antibody (AMA):} For PBC diagnosis.
  \item \textbf{p-ANCA:} Associated with PSC.
  \item \textbf{CA 19-9:} Tumor marker for pancreatic/biliary malignancy.
  \item \textbf{Liver biopsy/fibroscan:} For definitive diagnosis when non-invasive workup is inconclusive.
\end{itemize}

\section*{References}
\bibliographystyle{unsrtnat}
\bibliography{references}

\clearpage
\section*{Regulatory Status and Authorization Date}
This chapter describes a test category, not a specific commercial product. FDA, CDSCO, EU IVDR, and MHRA approval or registration dates therefore cannot be assigned without the assay manufacturer, product name, instrument, and intended use. For laboratory-developed tests, an individual agency approval date may not exist. See the Preface for the interpretation of regulatory status.

\clearpage
